% !PRINTED 486
\weberpage{486}
%<<CONT
If we then make, for the seven variables $\eta_r$, the substitution

\weq{4}{\eta_r = a_0 + a_1 \lambda_r + a_2 \lambda_r^2 + a_3 \lambda_r^3 + a_4 \lambda_r^4 + a_5 \lambda_r^5 + a_6 \lambda_r^6,}

whose determinant, being the product of all the differences $\lambda_i - \lambda_k$, is different from zero, and substitute correspondingly

\weq{5}{\eta_r = \sum_{0,\,6}^{s} a_s \lambda_r^s, \quad \eta_r' = \sum_{0,\,6}^{s} a_s' \lambda_r^s, \quad \eta_r'' = \sum_{0,\,6}^{s} a_s'' \lambda_r^s,}

then $X_1$ thereby goes over into a trilinear form of the three sequences of variables $a_s, a_s', a_s''$, which cannot vanish identically, since conversely the $a_s, a_s', a_s''$ are in turn linearly expressible in terms of the $\eta_r, \eta_r', \eta_r''$. Now one can assume for the variables $a_s, a_s', a_s''$ such rational numerical values (Bd. I, §.\,143) that the functions $X_1, X_2, X_3$ take values different from zero, and then (5) represents a suitable assumption for the functions $\eta_r, \eta_r', \eta_r''$\footnote{Cf., on the question hereby settled: \gesperrt{Burckhardt}, ``Ueber einen fundamentalen Satz der Lehre von den endlichen Gruppen linearer Substitutionen''. Mathem. Annalen, Bd.\ 42 (1892).}.
